\section{基于提示的推理方法}

\subsection{Chain-of-Thought 提示}

Chain-of-Thought（CoT）提示是激发语言模型推理能力的最基本方法。其核心思想是：让模型在给出最终答案之前，先产生\textbf{中间推理步骤}（rationale）。

CoT 提示有两种主要形式：

\begin{enumerate}
    \item \textbf{Few-shot CoT}：在提示中提供带有推理步骤的示例，模型在测试时模仿这些示例的推理模式。
    \item \textbf{Zero-shot CoT}：仅在提示末尾添加``Let's think step by step''，无需提供示例，模型即可产生推理过程。
\end{enumerate}

\begin{figure}[H]
\centering
\includegraphics[width=0.95\textwidth]{slides-images/slide-010.jpg}
\caption{Zero-shot CoT prompting 示例：(a) Few-shot 直接回答错误；(b) Few-shot CoT 提供推理步骤后正确；(c) Zero-shot 直接回答错误；(d) Zero-shot CoT 仅添加``Let's think step by step''即可正确推理\protect\footnotemark}
\end{figure}
\footnotetext{Slides 第11页。Source: Kojima et al. 2023。}

\begin{importantbox}{CoT 的核心洞察}
让模型``展示工作过程''（show its work）——即在输出答案前先生成推理步骤——可以显著提升模型在数学、逻辑等推理任务上的表现。这一发现表明，\textbf{中间计算步骤}对于复杂任务至关重要，即使这些步骤只是文本形式的。
\end{importantbox}

\subsection{Self-Consistency：多数投票}

Self-Consistency（自一致性）是对 CoT 的直接改进。标准 CoT 使用贪婪解码生成一条推理路径和答案；Self-Consistency 则\textbf{采样多条推理路径}，选择\textbf{出现频率最高的答案}。

\begin{figure}[H]
\centering
\includegraphics[width=0.95\textwidth]{slides-images/slide-012.jpg}
\caption{Self-Consistency 在多个数学推理 benchmark 上的表现：在所有模型和数据集上均显著优于标准 CoT prompting\protect\footnotemark}
\end{figure}
\footnotetext{Slides 第13页。Source: Wang et al. 2023。}

\begin{knowledgebox}{Self-Consistency vs 集成学习}
Self-Consistency 看似类似于传统的集成学习（Ensemble），但实验表明它的效果\textbf{优于}简单的 prompt 集成（同一模型用不同 prompt 后投票）。这可能是因为 Self-Consistency 通过采样不同的推理路径，探索了更多样化的解题思路，而非仅仅是对同一种推理的随机变体。
\end{knowledgebox}

\subsection{Least-to-Most Prompting：问题分解}

Least-to-Most Prompting 引入了\textbf{问题分解}策略：将大问题拆分为子问题，依次求解，最后综合得出最终答案。

\begin{figure}[H]
\centering
\includegraphics[width=0.95\textwidth]{slides-images/slide-015.jpg}
\caption{Least-to-Most Prompting 的两个阶段：Stage 1 将问题分解为子问题；Stage 2 依次求解子问题，后续子问题的回答以前序子问题的答案为条件\protect\footnotemark}
\end{figure}
\footnotetext{Slides 第16页。Source: Zhou et al. 2023。}

Least-to-Most Prompting 的一个有趣发现是：即使提示中仅包含需要\textbf{2步}推理的示例，模型也能泛化到需要\textbf{5步以上}推理的问题——这种泛化能力优于标准 CoT。但作者也指出，经过充分 prompt 工程的标准 CoT 可以达到类似的效果。

\begin{warningbox}{提示方法的局限性}
所有基于提示的方法都依赖于大语言模型（通常是 GPT-4 级别），且效果\textbf{高度依赖提示的设计}。对于较小的模型，这些提示方法往往效果有限。这引出了下一个问题：能否通过训练让小模型也具备推理能力？
\end{warningbox}

\subsection{本章小结}

基于提示的推理方法从简单到复杂包括：（1）Chain-of-Thought——让模型生成推理步骤；（2）Self-Consistency——采样多条路径后多数投票；（3）Least-to-Most Prompting——将问题分解为子问题逐步求解。这些方法都不修改模型参数，仅通过巧妙的提示设计来引导推理行为。

%% ========== Part 3: 训练方法 ==========

